\begin{table}[ht]
  \centering
  \caption{Policy Simulation: Percentage Change from Baseline at Period 10}
  \label{tab:policy_simulation}
  \small
  \begin{tabular}{lccc}
    \toprule
    Scenario & Outsourcing (\%) & Informal Wage (\%) & Employment (\%) \\
    \midrule
    Zero Enforcement ($E_s$ = 0) & -2.85 & 4.63 & 4.11 \\
    Double Enforcement ($E_s$ = 0.45) & 2.43 & -3.00 & -4.06 \\
    Informal Wage Floor (+20\%) & 0.26 & 7.12 & 0.00 \\
    Double $E_s$ + Wage Floor & 2.12 & 3.01 & -4.02 \\
    \bottomrule
  \end{tabular}
  \medskip
  \begin{minipage}{\linewidth}\footnotesize
    \textit{Notes:} 10,000 Monte Carlo draws from OLS sampling
    distributions (Stages 1--3). Baseline $E_s = 0.223$ (sample mean).
    Double enforcement $E_s = 0.446$. Wage floor: 20\% exogenous
    increase in informal wages from period 5.
    Persistence equation intercept anchored so baseline path is stable at
    observed mean ($\ln N = 4.44$, $N = 84.5$ enterprises).
    Wage $E_s$ coefficient ($\hat{\beta} = 0.114$, SE $= 5.09$)
    draws truncated at $\pm 3$, reflecting near-zero identification.
    The divergence between wage (Panel B) and employment (Panel C)
    responses to enforcement is the monopsony mechanism:
    more informal workers are drawn into the system as enforcement rises,
    but their wages remain suppressed.
  \end{minipage}
\end{table}
